\section{Hipótesis y justificación}
El problema principal es en el conflicto del incremento rápido del Big Data y entre el limitado de los recursos energéticos y su sostenibilidad en entornos de nube, aunque el cloud computing ha presentado soluciones de escalabilidad y seguridad [Artificial Intelligence Data Security Evaluation], pero las gestiones estos datos necesitan un consumo masivo de energía debido al funcionamiento continuo de los servidores y  los sistemas de la refrigeración, El problema de investigación es la ineficiencia de asignar dinámicamente los recursos en la nube que los recursos siguen trabajando al máximo de sus potencias incluso en periodos de baja demanda lo que resulta en desperdicio energético y un impacto ambiental. 

Este estudio se basa en hipótesis de la integración de modelos de predicción avanzados basados en Inteligencia Artificial (como el Deep learning) para analizar los patrones de flujo de dato para que los sistemas de computación puedan reasignar sus recursos dinámicamente y de forma proactiva y esto resulta en una reducción muy clara del consumo energético sin afectar al rendimiento de los servicios ni la seguridad de los datos.
